\subsection{Share price over the last 10 years}

\begin{figure}[H]
  \centering
  \includegraphics[width=0.62\textwidth]{z006_1.png}
  \caption{Boeing (BA) -- closing price over the last 10 years.}
\end{figure}

\begin{figure}[H]
  \centering
  \includegraphics[width=0.62\textwidth]{z008_1.png}
  \caption{Airbus (AIR.PA) -- closing price over the last 10 years.}
\end{figure}

\begin{figure}[H]
  \centering
  \includegraphics[width=0.62\textwidth]{z010_1.png}
  \caption{Toyota (TM) -- closing price over the last 10 years.}
\end{figure}

\paragraph{Comments.}
The price chart shows the time sequence on the abscissa and the price of the share on the
ordinate. Both Boeing and Airbus show a significant drop in price at the beginning of
2020, i.e.\ the first lockdown and the suspension of air traffic. While Airbus was almost
back to its previous high of 2020 at the beginning of 2022, Boeing shows a much more
negative trend. Toyota, on the other hand, reached its peak at the beginning of 2022, but
has also shown a negative trend since then.

\subsection{Price-earnings ratio (P/E ratio)}

The P/E values obtained for the three companies are summarised in Table~\ref{tab:pe}.

\begin{table}[H]
  \centering
  \caption{Price-earnings ratio (P/E).}
  \label{tab:pe}
  \begin{tabular}{lll}
    \toprule
    Company & P/E & Remark\\
    \midrule
    Boeing & -- & value not returned by the data source\\
    Airbus & 19.89 & \\
    Toyota & 10.37 & \\
    \bottomrule
  \end{tabular}
\end{table}

\paragraph{Comments.}
Given the P/E values for Airbus and Toyota of 19.9 and 10.4, respectively, it can be seen
that Airbus is above the guideline value of 15 and can thus be regarded as rather
expensive, while Toyota can be regarded as rather inexpensive by industry standards.

\subsection{Price-to-sales ratio (P/S ratio), last 12 months}

\begin{table}[H]
  \centering
  \caption{Price-to-sales ratio (P/S).}
  \label{tab:ps}
  \begin{tabular}{ll}
    \toprule
    Company & P/S\\
    \midrule
    Boeing & 1.502\\
    Airbus & 1.458\\
    Toyota & 0.006\\
    \bottomrule
  \end{tabular}
\end{table}

\paragraph{Comments.}
Similar to the P/S value above, it can be seen that both Boeing and Airbus are classified
as more or less expensive with values of 1.5 and 1.4 respectively, while Toyota could
even be classified as favorable with a value close to zero.

\subsection{Log-returns}

\begin{table}[H]
  \centering
  \caption{Cumulative log-returns over the observed period.}
  \label{tab:logret}
  \begin{tabular}{lrr}
    \toprule
    Company & Log-return & in \%\\
    \midrule
    Boeing & $-0.3025$ & $-30.2$\\
    Airbus & $-0.1451$ & $-14.5$\\
    Toyota & $-0.1486$ & $-14.9$\\
    \bottomrule
  \end{tabular}
\end{table}

\paragraph{Comments.}
The log return of the last year shows for all three stocks that a simple buy-and-hold
strategy would have led to more or less larger losses during this period. In the case of
Boeing with a loss of 30\%, in the case of Airbus with 14\% and for Toyota with a loss of
15\%.
